\begin{lemma}{2.1.8}\label{XXIII.2.1.8}
Let $n\in\uNorm_G(T)(S)$. For every $S'\to S$ and every $g\in\Omega(S')$ such that $\operatorname{int}(n)g\in\Omega(S')$, one has
\[
f(ngn^{-1})=f_N(n)f(g)f_N(n)^{-1}.
\]
\end{lemma}
\begin{proof}
This is trivial if $n\in T(S)$ (by 2.1.3). The two sides of the preceding formula define morphisms from $\Omega\cap\operatorname{int}(n)^{-1}\Omega$ to H; to check that they coincide, it suffices to check that there exists an open subset $V_n$ of $\Omega$, containing the identity section, such that $\operatorname{int}(n)V_n\subset\Omega$, and that $f\circ\operatorname{int}(n)$ and $\operatorname{int}(f_N(n))\circ f$ coincide on $V_n$. By the structure of $\uNorm_G(T)$, it suffices to prove that if the preceding assertion holds for an $n'\in\uNorm_G(T)(S)$ and if $\alpha\in\Delta$, it holds for $n=n'w_\alpha$. Put
\[
V_n=\Omega\cap\operatorname{int}(w_\alpha)^{-1}V_{n'}.
\]
One has $\operatorname{int}(n)V_n\subset\operatorname{int}(n')V_{n'}\subset\Omega$. If $g\in V_n(S')$, one has
\[
\operatorname{int}(n)g=\operatorname{int}(n')\operatorname{int}(w_\alpha)g.
\]
Now $\operatorname{int}(w_\alpha)g\in V_{n'}(S')$, so by hypothesis
\[
f\bigl(\operatorname{int}(n')\operatorname{int}(w_\alpha)g\bigr)
=\operatorname{int}(f_N(n'))f(\operatorname{int}(w_\alpha)g);
\]
since $\operatorname{int}(w_\alpha)g\in\Omega(S')$, one may apply 2.1.7, which gives
\[
f(\operatorname{int}(w_\alpha)g)=\operatorname{int}(f_N(w_\alpha))f(g),
\]
and one concludes immediately.
\end{proof}

Let us now prove 2.1.2. Let $x,x'\in\Omega(S')$; write as usual
\[
x=vtu,\qquad x'=v't'u',
\]
whence
\[
xx'=vt(uv')t'u'.
\]
By 2.1.6 and the relation $U^-(S')\Omega(S')U(S')=\Omega(S')$, one is reduced to proving that if $uv'\in\Omega(S')$, one has $f(uv')=f(u)f(v')$. Let $n\in\uNorm_G(T)(S')$ be as in 2.1.5 (ii).
% typo?: reference 2.1.5 (ii), kept as printed.
One has
\[
\begin{gathered}
f(u)=f_N(n)^{-1}f(nun^{-1})f_N(n),\\
f(v')=f_N(n)^{-1}f(nv'n^{-1})f_N(n),
\end{gathered}
\]
by \loccit, whence
\[
f(u)f(v')=f_N(n)^{-1}f(nun^{-1})f(nv'n^{-1})f_N(n).
\]
But $nun^{-1}\in U^-(S')$, $nv'n^{-1}\in U(S')$, so that
\[
f(nun^{-1})f(nv'n^{-1})=f\bigl((nun^{-1})(nv'n^{-1})\bigr)=f(nuv'n^{-1}),
\]
which gives
\[
f(u)f(v')=f_N(n)^{-1}f(nuv'n^{-1})f_N(n).
\]
If $uv'\in\Omega(S')$, one may apply 2.1.8 and the proof is complete.

\begin{remark}{2.2}\label{XXIII.2.2}
Instead of giving $f_N$, one may give a group morphism $f_T:T\to H$, and sections $h_\alpha\in H(S)$ $(\alpha\in\Delta)$ satisfying the conditions of 1.8.2. One must then replace condition (ii) of the theorem by

(ii$'$) There exists a group morphism $F_\alpha:Z_\alpha\to H$ which extends
\[
f_\alpha,\ f_{-\alpha},\ f_T\quad\text{and satisfies }F_\alpha(w_\alpha)=h_\alpha.
\]
\end{remark}

We shall now give the preceding theorem a more explicit form. Retain the preceding notation.
\begin{theorem}{2.3}\label{XXIII.2.3}
Let S be a scheme, and G a pinned S-group. Let H be an S-group sheaf for \fppf,
\[
f_T:T\to H,\qquad f_\alpha:U_\alpha\to H,\quad\alpha\in\Delta,
\]
group morphisms, and
\[
h_\alpha\in H(S),\qquad(\alpha\in\Delta),
\]
sections of H. For there to exist a group morphism
\[
f:G\to H
\]
(necessarily unique) inducing $f_T$ and the $f_\alpha$ $(\alpha\in\Delta)$ and satisfying $f(w_\alpha)=h_\alpha$ for every $\alpha\in\Delta$, it is necessary and sufficient that the following conditions hold:
\begin{enumeratei}
\item For every $S'\to S$, every $t\in T(S')$, every $\alpha\in\Delta$ and every $x\in U_\alpha(S')$, one has
\begin{equation}\tag{1}
f_T(t)f_\alpha(x)f_T(t)^{-1}=f_\alpha(\operatorname{int}(t)x).
\end{equation}
\item For every $\alpha\in\Delta$, every $S'\to S$, every $t\in T(S')$, one has
\begin{equation}\tag{2}
h_\alpha f_T(t)h_\alpha^{-1}=f_T(s_\alpha(t))=f_T(t\cdot\alpha^*\alpha(t)^{-1}).
\end{equation}
\item For every $(\alpha,\beta)\in\Delta\times\Delta$, one has
\begin{equation}\tag{3}
(h_\alpha h_\beta)^{n_{\alpha\beta}}=f_T(t_{\alpha\beta}).
\end{equation}
\item For every $\alpha\in\Delta$, one has (recall that one writes $u_\alpha=\exp_\alpha(X_\alpha)$)
\begin{equation}\tag{4}
(h_\alpha f_\alpha(u_\alpha))^3=e.
\end{equation}
\item For every $(\alpha,\beta)\in\Delta\times\Delta$, $\alpha\ne\beta$, there exists a group morphism
\[
f_{\alpha\beta}:U_{\alpha\beta}\to H
\]
satisfying the following two conditions:

 a) One has
\begin{equation}\tag{5}
f_{\alpha\beta}|_{U_\alpha}=f_\alpha,\qquad f_{\alpha\beta}|_{U_\beta}=f_\beta,
\end{equation}
 b) For every $\gamma\in R^+_{\alpha\beta}$, $\gamma\ne\alpha$ (resp. $\gamma\ne\beta$), and every $x\in U_\gamma(S')$, $S'\to S$, one has
% typo?: w_alpha in the beta alternative of (6), kept as printed.
\begin{equation}\tag{6}
\begin{gathered}
\operatorname{int}(h_\alpha)f_{\alpha\beta}(x)=f_{\alpha\beta}(\operatorname{int}(w_\alpha)x),\\
\text{(resp.}\quad\operatorname{int}(h_\beta)f_{\alpha\beta}(x)=f_{\alpha\beta}(\operatorname{int}(w_\alpha)x)\text{).}
\end{gathered}
\end{equation}
\end{enumeratei}
\end{theorem}
\emph{Proof.} Uniqueness is clear by 1.9.2. Let us prove existence.
\begin{lemma}{2.3.1}\label{XXIII.2.3.1}
There exists a group morphism
\[
f_N:\uNorm_G(T)\to H
\]
extending $f_T$ and satisfying $f_N(w_\alpha)=h_\alpha$.
\end{lemma}
Indeed, this is what 1.8.2 asserts, taking conditions (2) and (3) into account.
\begin{lemma}{2.3.2}\label{XXIII.2.3.2}
There exists a group morphism
\[
F_\alpha:Z_\alpha\to H,\qquad\alpha\in\Delta,
\]
extending $f_T$, $f_\alpha$ and satisfying $F_\alpha(w_\alpha)=h_\alpha$, hence extending $f_N|_{\uNorm_{Z_\alpha}(T)}$.
\end{lemma}
This is clear by Exp. XX, 6.2 and conditions (1), (3) and (4).
\begin{lemma}{2.3.3}\label{XXIII.2.3.3}
For every $(\alpha,\beta)\in\Delta\times\Delta$, $\alpha\ne\beta$, every $n\in\uNorm_{Z_{\alpha\beta}}(T)(S)$ such that $\bar n(\alpha)=\alpha$ (resp. $\bar n(\alpha)=\beta$), i.e., $\operatorname{int}(n)U_\alpha=U_\alpha$ (resp. $\operatorname{int}(n)U_\alpha=U_\beta$), every $S'\to S$ and every $x\in U_\alpha(S')$, one has
\[
\begin{aligned}
\operatorname{int}(f_N(n))f_\alpha(x)&=f_\alpha(\operatorname{int}(n)x)\\
\text{resp.}\quad\operatorname{int}(f_N(n))f_\alpha(x)&=f_\beta(\operatorname{int}(n)x).
\end{aligned}
\]
\end{lemma}
\begin{proof}
Indeed, there exist a $t\in T(S)$ and a sequence $\{\alpha_i\}\subset\{\alpha,\beta\}$ such that $n=tw_{\alpha_k}\cdots w_{\alpha_1}$. The condition holds for $n=t$ by condition (1). One may therefore suppose $n=w_{\alpha_k}\cdots w_{\alpha_1}$. Let us give the proof by induction on k, i.e., suppose the assertion proved for every n which can be written as a product of fewer than $k-1$ simple reflections (and satisfies the hypotheses of the lemma). Consider the roots
\[
\gamma_i=s_{\alpha_i}\cdots s_{\alpha_1}(\alpha).
\]
If all the $\gamma_i$ are positive, i.e., belong to $R^+_{\alpha\beta}$, one may apply condition (v) of 2.3; adopting the notation of 2.3 (v), one concludes immediately by induction that
\[
\operatorname{int}\bigl(f_N(w_{\alpha_i}\cdots w_{\alpha_1})\bigr)f_\alpha(x)
=f_{\alpha\beta}(\operatorname{int}(w_{\alpha_i}\cdots w_{\alpha_1})x),
\]
which, for $i=k$, is the desired property. If not all the $\gamma_i$ are positive, there exists an index $j<k$ such that
\[
\gamma_j\in R_+,\qquad\gamma_{j+1}\notin R_+.
\]
Since $\gamma_{j+1}=s_{\alpha_j}(\gamma_j)$, it follows immediately that $\gamma_j=\alpha_j$, by Exp. XXI, 3.3.2, hence that $\gamma_j$ is $\alpha$ or $\beta$, and one may decompose n as
% typo?: alpha_j rather than alpha_{j+1}, and fewer than k-1 in this proof, kept as printed.
\[
\begin{aligned}
n&=n'\cdot n'',\\
n'&=w_{\alpha_k}\cdots w_{\alpha_{j+1}},\\
n''&=w_{\alpha_j}\cdots w_{\alpha_1},
\end{aligned}
\]
$n'$ and $n''$ satisfying the hypotheses of the lemma and being products of fewer than $k-1$ reflections, hence satisfying the conclusion of the lemma by the induction hypothesis.
\end{proof}
\begin{lemma}{2.3.4}\label{XXIII.2.3.4}
Let $\alpha\in R$. If $n,n'\in\uNorm_G(T)(S)$ and $\beta,\beta'\in\Delta$ satisfy $\bar n(\alpha)=\beta$ and $\bar n'(\alpha)=\beta'$, one has
% typo?: f_T evaluated on n,n' in the normalizer, kept as printed.
\[
\operatorname{int}\bigl(f_T(n)^{-1}\bigr)f_\beta(nxn^{-1})
=\operatorname{int}\bigl(f_T(n')^{-1}\bigr)f_{\beta'}(n'xn'^{-1}),
\]
for every $x\in U_\alpha(S')$, $S'\to S$.
\end{lemma}
\begin{proof}
It suffices to check that if $\bar n(\alpha)=\beta$, $\alpha,\beta\in\Delta$, then $\operatorname{int}(f_T(n))\circ f_\alpha=f_\beta\circ\operatorname{int}(n)$. Now, by Tits's lemma (Exp. XXI, 5.6), there exist a sequence of simple roots $\alpha_0=\alpha$, $\alpha_1,\ldots,\alpha_m=\beta$, and a sequence of elements $w_i\in W$, $i=0,1,\ldots,m-1$, with
\[
\begin{aligned}
\bar n&=w_{m-1}\cdots w_0,\\
w_i(\alpha_i)&=\alpha_{i+1},\qquad i=0,1,\ldots,m-1,
\end{aligned}
\]
the following condition being moreover satisfied: for every $i=0,1,\ldots,m-1$, there exists a simple root $\beta_i$ such that
\[
w_i\in W_{\alpha_i\beta_i},\qquad\alpha_{i+1}=\alpha_i\quad\text{or}\quad\beta_i.
\]
One is then reduced by induction to the case treated in the preceding lemma.
\end{proof}
\begin{lemma}{2.3.5}\label{XXIII.2.3.5}
There exists a family of group morphisms $f'_\gamma:U_\gamma\to H$, $\gamma\in R$, satisfying
\begin{enumeratei}
\item For $\alpha\in\Delta$, one has $f'_\alpha=f_\alpha$ and $f'_{-\alpha}=F_\alpha|_{U_{-\alpha}}$, where $F_\alpha$ is defined in 2.3.2.
\item For $\alpha,\beta\in\Delta$ and $\gamma\in R^+_{\alpha\beta}$, one has $f'_\gamma=f_{\alpha\beta}|_{U_\gamma}$.
\item For every $n\in\uNorm_G(T)(S)$ and $\alpha,\beta\in R$ such that $\bar n(\alpha)=\beta$, one has
\[
\operatorname{int}(f_N(n))f'_\alpha(x)=f'_\beta(\operatorname{int}(n)x)
\]
for every $x\in U_\alpha(S')$, $S'\to S$.
\end{enumeratei}
\end{lemma}
\begin{proof}
For every root $\alpha\in R$, there exists an $n\in\uNorm_G(T)(S)$ such that $\bar n(\alpha)\in\Delta$. Define $f'_\alpha(x)$ to be the expression of 2.3.4. This is independent of the choice of n and $f'_\alpha$ is indeed a group morphism. Property (iii) holds by construction. The first assertion of (i) is clear (take $n=e$), and so is the second (take $n=w_\alpha$ and apply 2.3.2); if $\gamma\in R^+_{\alpha\beta}$, $(\alpha,\beta\in\Delta)$, there exists $n\in\uNorm_{Z_{\alpha\beta}}(T)(S)$ such that $\bar n(t)=\alpha$ or $\beta$; one then applies (iii) and conditions (5) and (6) and has proved (ii).
% typo?: n-bar(t) rather than n-bar(gamma), kept as printed.
\end{proof}
\numpara{2.3.6}
Let us now prove the theorem by proving that the conditions of 2.1 hold for the morphisms $f_N$ and $f'_\alpha$, $\alpha\in R$.

2.1 (i) follows from 2.3.5 (iii),

2.1 (ii) follows from 2.3.5 (i) and 2.3.2,

2.1 (iii) follows from 2.3.5 (ii) and the fact that $f_{\alpha\beta}$ is a group morphism.

An obvious corollary is:
\begin{corollary}{2.4}\label{XXIII.2.4}
Let S be a scheme, G a pinned S-group of semisimple rank $\geq1$, and H an \fppf S-group sheaf. For each $(\alpha,\beta)\in\Delta\times\Delta$, let
\[
F_{\alpha\beta}:Z_{\alpha\beta}\to H
\]
be a group morphism. For there to exist a group morphism
\[
f:G\to H
\]
inducing the $F_{\alpha\beta}$, it is necessary and sufficient that for every $(\alpha,\beta)\in\Delta\times\Delta$, one have
\[
F_{\alpha\beta}|_{Z_\alpha}=F_{\alpha\alpha}.
\]
\end{corollary}
\begin{proof}
The condition is obviously necessary. Suppose it satisfied. Put $f_T=F_{\alpha\alpha}|_T$ (which does not depend on $\alpha$, since $F_{\alpha\alpha}|_T=F_{\alpha\beta}|_T=F_{\beta\beta}|_T$). Put
% typo?: p_alpha instead of f_alpha, kept as printed.
\[
p_\alpha=F_{\alpha\alpha}|_{U_\alpha},\qquad h_\alpha=F_{\alpha\alpha}(w_\alpha),\qquad f_{\alpha\beta}=F_{\alpha\beta}|_{U_{\alpha\beta}}.
\]
The conditions of 2.3 are obviously satisfied: (1), (2), (4) ``are checked'' in $Z_\alpha$, (3), (5) and (6) in $Z_{\alpha\beta}$. There thus exists a morphism $f:G\to H$ extending $f_T$, the $f_\alpha$, $\alpha\in\Delta$, and satisfying $f(w_\alpha)=h_\alpha$; it coincides with $F_{\alpha\beta}$ on generators of $Z_{\alpha\beta}$, hence satisfies $f|_{Z_{\alpha\beta}}=F_{\alpha\beta}$.
\end{proof}
One also has the following technical corollary:
\begin{corollary}{2.5}\label{XXIII.2.5}
Let S be a scheme, G and $G'$ two pinned S-groups of semisimple rank 2, q an integer $>0$ such that $x\mto x^q$ is an endomorphism of $\mathbf G_{\mathrm a,S}$, and $h:\mathscr R(G')\to\mathscr R(G)$ a q-morphism of pinned root data. For each $\alpha\in R_+$, choose an $X_\alpha\in\Gamma(S,\goth g^\alpha)^\times$ and an $X'_{d(\alpha)}\in\Gamma(S,\goth g'^{d(\alpha)})^\times$ (extending the canonical choices for $\alpha\in\Delta$). Suppose the following conditions hold:
\begin{enumeratei}
\item If $\Delta=\{\alpha,\beta\}$, one has $D_S(h)(t_{\alpha\beta})=t'_{d(\alpha)d(\beta)}$.
\item For every $\alpha\in\Delta$ and every $\beta\in R_+$, $\beta\ne\alpha$ (whence $s_\alpha(\beta)\in R_+$), if $z\in\Gm(S)$ is defined by
\[
\Ad(w_\alpha)X_\beta=zX_{s_\alpha(\beta)},
\]
one also has
\[
\Ad(w'_{d(\alpha)})X'_{d(\beta)}=z^{q(\beta)}X'_{d(s_\alpha(\beta))}.
\]
\item There exists a group morphism $f:U\to U'$ such that for every $\alpha\in R_+$ one has, for every $x\in\Ga(S')$, $S'\to S$,
\[
f(\exp(xX_\alpha))=\exp(x^{q(\alpha)}X'_{d(\alpha)}).
\]
\end{enumeratei}
Then there exists a morphism of pinned groups $G\xrightarrow{g}G'$ such that $\mathscr R(g)=h$.
\end{corollary}
\begin{proof}
Indeed, one defines $f_\alpha:U_\alpha\to G'$ by
\[
f_\alpha(\exp(xX_\alpha))=\exp(x^{q(\alpha)}X'_{d(\alpha)});
\]
one puts $f_T=D_S(h)$, $h_\alpha=w'_{d(\alpha)}$. The conditions of 2.3 are satisfied (observe that $q(s_\alpha(\beta))=q(\beta)$ (Exp. XXI, 6.8.4) and that one always has $D_S(h)(t_\alpha)=t'_{d(\alpha)}$) and one concludes immediately.
\end{proof}
\begin{remark}{2.6}\label{XXIII.2.6}
One may make condition (v) of 2.3 more precise as follows. One must first check:

(a) For every word in $w_\alpha$ and $w_\beta$ such that the corresponding word takes $\alpha$ or $\beta$ to $\alpha$ or $\beta$, the corresponding relation of type 2.3.3 holds. In fact the proof of 2.3.3 shows that it suffices to check it for the words in $w_\alpha$ and $w_\beta$ which are minimal (in the sense that every nontrivial initial subword fails to satisfy the imposed conditions).

If condition (a) holds, one may now define, for each $\gamma\in R^+_{\alpha\beta}$, an $f_\gamma:U_\gamma\to H$ as in 2.3.5; one must then write:

(b) The morphism defined by the $f_\gamma$
\[
U_{\alpha\beta}=\prod_{\gamma\in R^+_{\alpha\beta}}U_\gamma\to H
\]
is a group morphism. According to Exp. XXII, 5.5.8, (b) is implied by:

(b$'$) The preceding morphism respects the commutation relations between $U_\gamma$ and $U_\delta$ for $\gamma,\delta\in R^+_{\alpha\beta}$, $\gamma>\delta$ (i.e., the relations in $C_{i,j,\gamma,\delta}$ of Exp. XII, 5.5.2).
% typo?: Exp. XII rather than XXII, kept as printed.

It is clear that, conversely, the collection of conditions (a) and (b$'$) is equivalent to (v).

One may even reduce the preceding conditions to conditions involving only the elements $h_\alpha$, $h_\beta$, $f_\alpha(u_\alpha)$, $f_\beta(u_\beta)$ of H. A condition of type (a) will be written, for example, if $s_\alpha s_\beta s_\alpha(\beta)=\alpha$:
\begin{equation}\tag{1}
\operatorname{int}(h_\alpha h_\beta h_\alpha)f_\beta(x)=f_\alpha(\operatorname{int}(w_\alpha w_\beta w_\alpha)x);
\end{equation}
for every $x\in U_\alpha(S')$, $S'\to S$. In particular, for $x=u_\beta$, one has $\operatorname{int}(w_\alpha w_\beta w_\alpha)u_\beta=u_\alpha^z$ for a certain section z of $\Gm(S)$, and the preceding relation gives
% typo?: x in U_alpha rather than U_beta in this argument, kept as printed.
\begin{equation}\tag{1$'$}
\operatorname{int}(h_\alpha h_\beta h_\alpha)f_\beta(u_\beta)=f_\alpha(u_\alpha)^z.
\end{equation}
Let us show that, conversely, taking conditions (i) to (iv) of 2.3 into account, (1$'$) implies (1). If $t\in T(S')$, $S'\to S$, let $\operatorname{int}(t)$ act on (1$'$); taking conditions (i) and (ii) into account, one obtains (1) for $x=\operatorname{int}(t)u_\beta=u_\beta^{\beta(t)}$. It now suffices to observe that $\beta:T\to\mathbf G_{\mathrm m,S}$ is faithfully flat, hence that condition (1) certainly holds for $x\in U_\alpha(S')^\times$, $S'\to S$. Since it is additive in x and every section of $U_\alpha$ can locally be written as a sum of two sections of $U_\alpha^\times$, one does indeed conclude that (1$'$) + (i) + (ii) $\Longrightarrow$ (1).

One argues in the same way with conditions of type (b). One must then use the fact that if $\gamma$ and $\gamma'$ are two distinct positive roots (and hence linearly independent over $\QQ$), the morphism $T\to\mathbf G_{\mathrm m,S}^2$ with components $\gamma$ and $\gamma'$ is faithfully flat. We leave the details of this transposition to the reader.
\end{remark}

\sgasection{3}{Groups of semisimple rank 2}
\numpara{3.1}\textbf{Generalities}
\begin{lemma}{3.1.1}\label{XXIII.3.1.1}
Let S be a scheme, G a pinned S-group, and $\alpha$ and $\beta$ two roots of G, with $\alpha+\beta\ne0$.
\begin{enumeratei}
\item If $\alpha+\beta\notin R$, one has
\[
\exp(X_\alpha)\exp(X_\beta)=\exp(X_\beta)\exp(X_\alpha)
\]
for all $X_\alpha\in\mathbf W(\goth g^\alpha)(S')$, $X_\beta\in\mathbf W(\goth g^\beta)(S')$, $S'\to S$.
\item If $\alpha+\beta$ and $\alpha-\beta$ are not roots, one has
\[
w_\alpha(z_\alpha)\exp(X_\beta)w_\alpha(z_\alpha)^{-1}=\exp(X_\beta)
\]
for all $X_\beta\in\mathbf W(\goth g^\beta)(S')$, $z_\alpha\in\mathbf W(\goth g^\alpha)^\times(S')$, $S'\to S$, and
\[
w_\alpha(z_\alpha)w_\beta(z_\beta)=w_\beta(z_\beta)w_\alpha(z_\alpha)
\]
for all $z_\alpha\in\mathbf W(\goth g^\alpha)^\times(S')$ and $z_\beta\in\mathbf W(\goth g^\beta)^\times(S')$, $S'\to S$.
